\documentclass[11pt,a4paper]{article}
\usepackage[margin=25mm]{geometry}
\usepackage[T1]{fontenc}
\usepackage{lmodern,microtype,amsmath,amssymb,amsthm,booktabs,array}
\usepackage[dvipsnames]{xcolor}
\usepackage[colorlinks=true,linkcolor=MidnightBlue,urlcolor=MidnightBlue,citecolor=MidnightBlue]{hyperref}
\usepackage{enumitem}
\setlist{itemsep=3pt,topsep=5pt}
\newtheorem{theorem}{Theorem}[section]
\newtheorem{lemma}[theorem]{Lemma}
\newtheorem{corollary}[theorem]{Corollary}
\newtheorem{proposition}[theorem]{Proposition}
\theoremstyle{remark}\newtheorem{remark}[theorem]{Remark}
\newcommand{\F}{\mathbb F}\newcommand{\Q}{\mathbb Q}\newcommand{\Z}{\mathbb Z}
\newcommand{\bits}{\{0,1\}}\newcommand{\poly}{\operatorname{poly}}
\title{Constructive phase representations\\\large Exact counting beyond the number of intermediate states}
\author{Astra P3 research}\date{27 October 2026}
\begin{document}\maketitle
\begin{abstract}
Conditional on the pinned OpenAI counting manuscript, we construct an exact
counter for products of fixed-locality, fixed $2$-power root-of-unity phases
on $n$ independent bits, with $m=n\exp(o(\log n))$ factors, in
$\poly(L+n)2^{o(n)}$ bit time and polynomial space. The interaction pattern
is unrestricted. A common divisibility filtration makes the retained input
degree $O(n)$ independently of $m$, while retaining one small scope per
local operation. The source's replacement forms and shared summaries are
then manufactured at growing incidence, rather than supplied.

The result gives exact preimage counting for checked bounded-degree
arithmetic programs whose compiled outputs meet the stated support and
output-count bounds, even with many intermediate states, and exact amplitudes
and requested marginals for circuits with $O(n)$ Hadamards and
$n^{1+o(1)}$ local diagonal phase gates. Raw descriptions are compiled,
with expansion, reconstruction, and arithmetic costs charged. Failed
repairs of the general-circuit problem are recorded as well. This does
not establish subexponential counting for arbitrary circuits or a
practical speed advantage. Executed programs validate the representation
and finite identities; they do not implement the source's full
asymptotic replacement solver.
\end{abstract}

\section{Result, dependency, and the obstruction removed}
All algorithmic claims using replacement forms are conditional on
\cite{governing}, specifically Definition 3.1, Proposition 3.4,
Lemma 4.7, and the rational arithmetic argument of Appendix B. The supplied
PDF was compared byte-for-byte with the pinned Git blob. The relevant
sparse-equation reduction, sampling and inversion consequences,
linear-size circuit consequence, and tensor-contraction appendix were
read before formulating this extension. This report does not claim an
independent verification of every theorem in the manuscript.

A direct circuit encoding counts on the original input bits plus gate
bits. Many intermediate states can erase the improvement measured in
the original inputs. Treating a global parity polynomial as one factor
avoids gate bits but gives it a large scope. Splitting it into local
monomial factors keeps scopes small, but ordinarily charges all factors
to the interpolation degree. We remove this last cost using divisibility
and pay explicitly for growing incidence on the original variables.

\begin{theorem}[Local phase products]\label{thm:main}
Fix integers $k,b\ge1$. Set $q=2^b$, $d=2^{b-1}$, and let $\zeta$ be a
primitive $q$-th root of unity. An input explicitly specifies $n$ Boolean
variables and $m$ tables $f_j:\bits^{I_j}\to\Z/q\Z$, with $|I_j|\le k$.
Let $L$ include all indices and tables. The exact algebraic integer
\[
 Z=\sum_{x\in\bits^n}\zeta^{\sum_{j=1}^m f_j(x|_{I_j})}
\]
can be computed deterministically in $\poly_{k,b}(L+n)2^{o(n)}$ bit
operations and polynomial space whenever
$\log(2+m/n)=o(\log n)$.
No graph-width, geometry, low-rank, or supplied-decomposition hypothesis
is imposed.
\end{theorem}

A polynomial exponent may depend on fixed $k,b$, but not on input
length or density. For $n\ge4$, set
\begin{align*}
 \lambda&=\lceil\log_2n\rceil,&\nu&=\lceil\log_2\lambda\rceil,\\
 h&=\lceil\log_2(2+km/n)\rceil+4\nu+\lceil\log_2(k+2)\rceil+10,
 &a&=h/\lambda.
\end{align*}
For sufficiently small $a$ the quantitative bound is
\begin{equation}\label{eq:time}
 (L+n)^{O_{k,b}(1)}
 2^{O_{k,b}\left(n\sqrt a\,[\log_2(2/a)]^{3/2}\right)}.
\end{equation}
One guarded algorithm handles all inputs, using enumeration outside the
useful range; Section~\ref{sec:guards} makes this explicit.

\paragraph{A strict regime beyond direct gate encoding.}
For $0<\alpha<1$, let $m=n2^{(\log_2n)^\alpha}$. The exponent becomes
\[
 O_{k,b,\alpha}\!\left(
 n\frac{(\log\log n)^{3/2}}{(\log n)^{(1-\alpha)/2}}\right)=o(n).
\]
A direct bounded-fan-in accumulation of $m$ local terms has
$\Theta(m)$ gate states. Substituting that variable count into the
governing circuit-to-CNF bound gives no improvement over $2^n$ in this
regime. This compares the explicit reductions, not every conceivable
encoding. The new bound removes a state-count and degree obstruction;
it is not just the linear-size circuit corollary. Single quadratic
$\F_2$ phase sums already have polynomial-time algorithms
\cite{williams2018}; the substantive binary application begins at degree
three. Multiple quadratic outputs are also relevant.

\section{A degree budget obtained from divisibility}
For $p(x)=\bigoplus_{S\in\mathcal M}x^S$ over $\F_2$, with $|S|\le k$,
form
\begin{equation}\label{eq:binaryproduct}
 F(X,V)=\prod_{S\in\mathcal M}(1-2VX^S)
       \quad\text{in }\Q[V]/(V^{n+2}).
\end{equation}
$V$ is a formal bookkeeping variable, not a Boolean state. Give each
factor independent incidence indeterminates before measuring degree.
A retained term of $V$-degree $j$ selects $j$ nonconstant factor parts,
so its incidence degree is at most $kj$. The full retained product has
degree at most $k(n+1)$, independently of $m$. This satisfies the exact
degree hypothesis of Proposition 3.4 without globally identifying and
multilinearizing all copies.

On Boolean inputs the untruncated product at $V=1$ is $(-1)^{p(x)}$.
Discarded orders $j\ge n+2$ are divisible by $2^{n+2}$. After reconstructing
the retained polynomial at the undeformed point,
\[
 \sum_x\sum_{j=0}^{n+1}[V^j]F(x,V)
 \equiv Z_p:=\sum_x(-1)^{p(x)}\pmod {2^{n+2}}.
\]
Since $Z_p\in[-2^n,2^n]$, centered reconstruction recovers it and
$\#\{x:p(x)=0\}=(2^n+Z_p)/2$. Just $n+1$ modular bits would identify
the endpoints $+2^n$ and $-2^n$.

For general phase tables, multilinearly interpolate
$x\mapsto\zeta^{f_j(x)}$ to $G_j\in\Z[\zeta][X_{I_j}]$, of degree
at most $k$ and with at most $2^k$ coefficients. This costs $O_k(1)$
operations per explicit table. Put $\pi=\zeta-1$. Modulo $(\pi)$ each
table is identically one. Boolean finite differences therefore show
that every coefficient of $G_j-1$ is divisible by $\pi$.

\begin{lemma}[Exact finite filtration]\label{lem:filtration}
Set $P=n+2$, $\ell=dP$, and
\[
 \mathcal A=\Q(\zeta)[V]/(V^\ell),\qquad B_j=1+V(G_j-1).
\]
The retained product $\prod_jB_j$ has incidence degree at most
$K=k(\ell-1)$. Its coefficient sum on Boolean inputs agrees with the
original phase product modulo $2^P\Z[\zeta]$.
\end{lemma}
\begin{proof}
Count the selected nonconstant parts for the degree bound.
Also $\pi^d\in2\Z[\zeta]$. For $d=1$, $\pi=-2$.
For $d\ge2$, expand $(\zeta-1)^d$: the endpoints cancel since
$\zeta^d=-1$, and all intermediate binomial coefficients are even
because $d$ is a power of two. An omitted term has at least $dP$
factors in $(\pi)$ and hence belongs to $2^P\Z[\zeta]$.
\end{proof}

Use the integral basis $1,\zeta,\ldots,\zeta^{d-1}$, with
$\zeta^d=-1$. Each input contributes one signed basis vector, so every
coordinate of $Z$ lies in $[-2^n,2^n]$. Its residues modulo $2^{n+2}$
determine all coordinates exactly. No ramification unit is needed.

\begin{remark}[A final linear map, not a homomorphism]
Setting $V=1$ is not an algebra homomorphism on the truncated ring.
Sum coefficients in its unique representative of degree below $\ell$
only after exact reconstruction at the undeformed point. The lemma
justifies this final map. The modular congruence is required only at
original Boolean inputs; the non-Boolean, rational replacements at
other points are handled by exact rational arithmetic.
\end{remark}

\section{Manufacturing the replacement forms}\label{sec:guards}
Add a neutral factor of scope $\{v\}$ for each variable, so total
incidence $M\le km+n$. Branch on variables belonging to more than
\[
 B_0=\lceil M/n\rceil\lceil\sqrt\lambda\rceil
\]
factor scopes. There are at most $n/\sqrt\lambda$ such variables.
Branches partition the input cube and leave $N\ge n/2$ free bits
for sufficiently large $n$. Substitute fixed bits in the local tables
and retain neutral factors. The maximum remaining incidence $B$ is at
most $B_0+1$, the scope bound $w$ is at most $\max(k,1)$, and
$\log_2B=O_k(h)$.

Put $\lambda'=\lceil\log_2N\rceil$ and recompute a comparable $h'$ from
the actual branch, retaining original precision $P=n+2$. Choose
\begin{align*}
 r&=\lfloor\sqrt{\lambda'/h'}\rfloor,&
 t&=\lceil\log_2r\rceil,&Q&=\lfloor\lambda'/(128t)\rfloor,\\
 \tau&=(r\lambda')^{-1},&D&=r^{2Q},&
 A&=\lceil\log_2(2Br^2\lambda')\rceil,\\
 R&=\lceil\sqrt{8QA}\rceil+32.
\end{align*}
Check $r\ge2$, $R\ge32$, $6R\le Q$, central-interval endpoints, and
the source's two balance conditions
\[
 2(B+w)D^2\le M_{\rm branch},\qquad 2Qr^2\le N.
\]
As $a\to0$,
\[
 A=O_k(h),\quad D\le(2N)^{1/64},\quad B=N^{o(1)},\quad
 R/Q=O_k(\sqrt{a\log(2/a)}).
\]
The balance guards thus hold eventually. An $R$-point central interval
contains the integers within $R/4$ of the source binomial mean. Its
exceptional fraction is bounded by
\[
 2\exp(-R^2/(8Q))\le 2^{1-A}
 \le (Br^2\lambda')^{-1}=\tau/(rB).
\]
Lemma 4.7 now constructs the occurrence-specific replacement rows,
partitions, and shared integer measurement matrices. Rows restore
original coordinates at zero and have degree at most $R$. Groups have
size at most $2N/r$. There are at most $2\tau N$ measurement rows in
total, each with absolute coefficient sum at most $D$. All are
manufactured in polynomial time in their explicit dimensions.

\paragraph{Uniform fallback.}
$R/Q$ need not tend to zero at a fixed positive $a$; this cannot be
hidden in a bounded-exceptions statement. Below a sufficiently small
fixed threshold for $a$, the guards hold uniformly for sufficiently
large $n$. Above it use original-input enumeration. Its $2^n$ cost is
absorbed by increasing the constant in \eqref{eq:time} on a bounded
positive-$a$ interval; larger $a$ simply uses the stated fallback.
Thus one algorithm has $2^{o(n)}$ cost along every family with $a\to0$.

\section{Summation, reconstruction, and bit costs}
Apply Proposition 3.4 in the commutative rational algebra $\mathcal A$,
of dimension $d\ell=O_b(n)$. The degree $K$ is $O_{k,b}(N)$ since
$N\ge n/2$. The Taylor correction order is
\[
 q_T=\lfloor KR/Q\rfloor=O_{k,b}(n\sqrt{a\log(2/a)}).
\]
For $q_T\le N/2$, the subset-pair enumeration satisfies
\[
 C_{q_T}=\sum_{j\le q_T}2^j\binom Nj,
 \qquad \log_2 C_{q_T}\le N(H_2(q_T/N)+q_T/N).
\]
This supplies the entropy logarithm in \eqref{eq:time}.
Measurement and group enumeration cost
\[
 (2D+1)^{2\tau N}2^{2N/r}=2^{O(N/r)}=2^{O(n\sqrt a)}.
\]
For other correction orders use $C_{q_T}\le3^N$, absorbed where the
claimed exponent is already a constant multiple of $n$. The high-incidence
branches add $n/\sqrt\lambda$, also absorbed.

Let $\delta=(R-1)!Q!$ be the source's row denominator. A retained
$V^j$ coordinate uses at most $kj$ replacement inputs, so its denominator
divides $\delta^{kj}$, including in Taylor evaluation. After clearing
denominators a coefficient-norm bound is a sum of terms of the form
\[
 \binom mj(C_{k,b}\max(\delta,B_*)^k)^j,
 \quad j<\ell,\qquad
 \log B_*=O(Q\log(Q+2)+\log(nD+2)).
\]
Its logarithm is polynomial in $L+n$ and the design parameters.
Summing exponential enumerations adds their logarithm to bit length.
Cyclotomic multiplication uses the fixed relation $\zeta^d=-1$.
Truncated Taylor/$V$ arrays and Hermite systems have polynomial
dimension, and rational reconstruction has polynomial bit cost by
the Appendix B argument coordinatewise. Reconstruct at $T=0$, sum
retained $V$ coefficients, reduce integral coordinates modulo $2^{n+2}$,
and center. Sequential traversal gives polynomial space. This proves
Theorem~\ref{thm:main}.

\section{Compiling raw arithmetic programs and preserving outputs}
Consider an explicitly supplied program over $\F_2$ or $\Z/2^b\Z$ with
$n$ Boolean inputs and $S$ arithmetic gates. For fixed formal degree $k$,
propagate canonical Boolean multilinear polynomials through the actual
gate DAG. Addition combines coefficient dictionaries; multiplication
takes pairwise unions of supports, using $x_i^2=x_i$, and combines
coefficients modulo the base ring. Each stored polynomial has at most
$\sum_{j\le k}\binom nj=O_k(n^k)$ terms. A conservative expansion cost is
\[
 O_{k,b}(S n^{2k})
\]
coefficient operations, with index-bit costs polynomial in the explicit
description. Formal degree is checked syntactically before claiming this
bound. The supplied Boolean compiler also accepts the broader class where
every intermediate canonical polynomial has degree at most $k$: it fully
forms each product of two degree-$k$ dictionaries, then checks the actual
degree before retaining it. A degree violation is reported, not silently
truncated. This uses the same conservative polynomial bound.

The final support is manufactured by this computation. An arbitrary
program merely promised to have a sparse final expansion is not assumed
to have an efficient compiler. If the constructed support is too large,
the favorable counting bound is not claimed. The result therefore
provides a checkable representation compiler for a restricted algebraic
class, rather than a universal compression assertion.

For a single output modulo $q$, compute all $q$ character sums and then
Fourier-invert after their exact reconstruction. One primitive character
alone need not determine a $q$-value histogram when $q>2$.

\begin{corollary}[Vector target-preimage counting]\label{cor:program}
Suppose the compiler obtains $t\le Cn$ output polynomials $P_i$ of degree
at most fixed $k$, with total nonzero support $m=n\exp(o(\log n))$.
For a supplied target $y$, its exact original-input preimage count is
computable in compilation cost plus $\poly(L+n)2^{o(n)}$ time and
polynomial space. There is no charge of one counting variable per
intermediate program state.
\end{corollary}
\begin{proof}
Introduce $bt$ character bits encoding $\lambda_i\in\Z/q\Z$. Orthogonality
gives the exact integer identity
\[
 \sum_{x,\lambda}\zeta^{\sum_{i=1}^t\lambda_i(P_i(x)-y_i)}
 =q^t\,\#\{x:P_i(x)=y_i\text{ for all }i\}.
\]
Expand $\lambda_i$ into its $b$ bits. The phase now has
$N=n+bt=O(n)$ independent bits, locality at most $k+1$, and at most
$bm+bt$ monomial gates. Theorem~\ref{thm:main} applies. Precision is
$N+2$, not $n+2$, so the multiplicity is fully charged. Reconstruct
the character sum exactly, then divide by $q^t$ exactly.
\end{proof}

For a raw local accumulation program the local factors are already explicit
instructions, so this polynomial-expansion stage is unnecessary. Its
arbitrarily long sequence of prefix states is replaced directly by the
product construction. In both formulations the resulting global output
may depend on all $n$ inputs; no constant output-locality hypothesis is used.

\section{Exact quantum amplitudes and requested marginals}
\subsection{IQP circuits}
For $U=H^{\otimes n}DH^{\otimes n}$, with $D$ a product of the permitted
local diagonal phase gates, the amplitude at $y$ is
\[
 \langle y|U|0^n\rangle=2^{-n}\sum_x\zeta^{P(x)}(-1)^{x\cdot y}.
\]
The extra phases are unary. For output positions $A$, $|A|=a_0$, and
$B=[n]\setminus A$, summing the unspecified output characters gives
\begin{equation}\label{eq:prefix}
 \Pr[Y_A=z]=2^{-(n+a_0)}
 \sum_{x_B,x_A,x'_A}
 \zeta^{P(x_A,x_B)-P(x'_A,x_B)}
 (-1)^{z\cdot(x_A+x'_A)}.
\end{equation}
Indeed orthogonality forces the two $B$-input copies to agree. This is
a phase sum on at most $2n$ bits with at most $2m+2a_0$ local factors,
so Theorem~\ref{thm:main} computes it exactly. No list of all $2^n$
amplitudes is produced for free.

\subsection{Arbitrary alternation with a linear Hadamard count}
The following extension was constructed and checked separately from
the IQP case. The path expansion itself is established machinery
\cite{decolnet}; the new complexity bound comes from the phase counter.

\begin{corollary}\label{cor:hadamard}
Fix $k,b,C$. Let a raw $n$-qubit circuit, on computational input $|0^n\rangle$,
have $h_H\le Cn$ Hadamard gates and $m=n\exp(o(\log n))$ diagonal gates
of locality at most $k$, with phases in the $2^b$-th roots of unity.
The gates may alternate arbitrarily and have arbitrary interaction
geometry. Every requested amplitude and every requested computational
output marginal is exactly computable in $\poly(L+n)2^{o(n)}$ bit time
and polynomial space, conditional on the governing manuscript.
\end{corollary}
\begin{proof}
Track the computational bit on each wire. A diagonal gate leaves this
bit unchanged and contributes a local phase on the current bits. At a
Hadamard with old bit $u$, introduce one new path bit $v$ and use
\[
 \langle v|H|u\rangle=2^{-1/2}(-1)^{uv}.
\]
Thus at most $h_H$ path bits are introduced, with no new bit at any
diagonal gate. Initial wire values are constants. The accumulated phase
has at most $m+h_H$ local terms of locality at most $\max(k,2)$.

For an amplitude, pin the final wire values to the requested output.
These are pins on path bits, or tests of constants. Contradictions give
zero. Eliminate consistent pins by substitution. Each remaining path
assignment corresponds to exactly one compatible computational history.
The amplitude is $2^{-h_H/2}$ times its phase sum.

For a marginal, take a path expression and its complex-conjugate copy.
Pin measured final wires in both copies. Identify the two final bits
on each unmeasured wire, implementing the sum over the unspecified output.
These equalities and pins are constructed by union-find. Their elimination
does not enlarge any local scope or introduce multiplicity. There are at
most $2h_H$ free variables and $2(m+h_H)$ local terms. The probability
is $2^{-h_H}$ times the resulting phase sum.

To keep a uniform parameter comparable to the physical qubit count,
pad either expression with unused Boolean variables to
$N_*=n+2h_H=\Theta(n)$. This multiplies the sum by a known power of two.
Use precision $N_*+2$, reconstruct exactly, and divide by this known
power after reconstruction. The factor density is $N_*^{o(1)}$.
Theorem~\ref{thm:main} applies. Return amplitudes as exact cyclotomic
coordinates times $2^{-h_H/2}$; if necessary $\sqrt2$ is adjoined to
the fixed number field. All normalization and boundary processing
have polynomial bit cost.
\end{proof}

The compiler and an independent state-vector evaluator agree exactly
on the finite tests in Section~\ref{sec:executed}. The restriction on
Hadamard count is substantive. Arbitrary gate sets cannot be converted
without charging any added Hadamards. Exact sampling is a standard
conditional-probability consequence in this extended family: fixed-degree
algebraic probabilities admit certified dyadic comparisons, and a random
dyadic threshold is resolved in expected polynomial work. No fast sampler
or practical simulation advantage is claimed to have been executed here.

\section{Why the direct general-circuit attacks failed}
The research did not stop at the first special case. The following
repairs were constructed, tested, and rejected at their precise scope.
Full proofs and exact checkers are included in the accompanying archive.

\subsection{Consistency tests do not eliminate arbitrary trace entropy}
For a topologically ordered Boolean circuit with inputs $x$ and stored
gate bits $y_1,\ldots,y_s$, define residuals
\[
 e_g=y_g+f_g(x,y_1,\ldots,y_{g-1})\quad\text{over }\F_2.
\]
The map $(x,y)\mapsto(x,e)$ is bijective: recover $y_g$ successively.
Even at fixed $x$, every one of the $2^s$ residual vectors occurs.
A rank-$r$ parity filter on residuals therefore leaves exactly
$2^{n+s-r}$ traces. Unlike the source's clause-violation vectors on
the original input cube, these residuals cannot all be removed by a
shorter exact linear consistency filter.

This is not a lower bound against output-specific algorithms. We tested
that escape too. Write $\chi(e)=(-1)^{e_1+\cdots+e_s}$. If a rational
weight $w$ preserves both the constant observable and the pure parity
observable at $e=0$, then $\sum_e w(e)\chi(e)=1$. Every polynomial
of degree less than $s$ is orthogonal to $\chi$, since each monomial
omits a coordinate. Thus such a selector has degree at least $s$, even
with signed weights. Alternatively a rank-one parity filter can preserve
this particular output, but its multiplicity is $2^{s-1}$. More generally
an invariant rank-$r$ filter gives $Z_{\rm filt}=2^{s-r}Z$; recovering
arbitrary $0\le Z\le2^n$ from a power-of-two residue requires
at least $n+s-r+1$ bits before division. A small filter alone does not
remove the state cost. The new output-character construction avoids this
failure by using only $O(n)$ character bits and charging their precision.

\subsection{A shared linear bottleneck has a degree budget}
If an $n\times n$ polynomial matrix $P(T)$ has entry degree at most $R$
and $P(0)=I$, then at distinct nodes $t_1,\ldots,t_Q$,
\[
 \sum_j\bigl(n-\operatorname{rank}P(t_j)\bigr)\le nR.
\]
Indeed, the rank deficiency at $t_j$ is a lower bound on the order of
vanishing of $\det P$ there, while $\det P$ is nonzero and has degree
at most $nR$. In particular, a common rank-$o(n)$ image at every cheap
node forces $R/Q\ge1-o(1)$. Incidence-specific forms are essential;
one cannot simply replace them by one small universal input summary.

The stronger signed-alias repair also failed for the tested complete
auxiliary grid. For $n=r^Q$, the source operator $M=Vf(W)$ with
$\deg f\le R-1$, and every nonzero rational signing/scaling $J_w$ of
replicated inputs, a valid summary satisfies
\[
 \operatorname{rank}(CJ_w)\ge
 \sum_{j=R}^{Q}\binom Qj(r-1)^j.
\]
At $R/Q=o(1)$ this is almost all original-input rank. The proof uses
the first exceptional spectral level and triangular equations on swap
cubes; the unsigned bound is exact. This rules out that repair of the
source grid, not all possible representations.

\subsection{Causality, precision, and the tensor bridge}
Imposing triangularity on the source replacement operator has an immediate
cost: every forced zero diagonal entry of $M$ contributes a pivot to
$I-M$. If there are $b_*$ such entries, any summary retaining its fixed
space has rank at least $b_*$. In a finite $64$-coordinate fixture the
noncausal source operator has defect rank $10$, whereas any triangular
operator satisfying the same support restriction has defect rank at
least $32$. Repair by deleting variables to make the auxiliary dependency
graph acyclic also consumes almost all variables in the tested regime.

The Appendix D interpolation precision cannot simply be replaced by the
output bit length. Explicit integer polynomial families have value one
at the target but vanish to precision $\Omega(MD/u)$ at all active
local evaluation nodes. A separate pinned-wire fixture exhibits the
same issue at small size. The present phase construction instead
truncates a globally justified polynomial \emph{before} rational
interpolation; it does not assume that tensor precision floor disappears.

N1077 was tested at its actual generic planar trivalent rational scope
\cite{n1077}. The tensor proof and packet's gauge admissibility conditions
were read. An AND-chain edge has zero cross-pairings in the relevant
canonical decomposition, and a separate nondegenerate rational theta
network fails the required compatibility products. A simple low-degree
deformation cannot supply arbitrarily many compatible nodes without
paying the cleared compatibility degree. These checks do not refute
N1077; they show why it does not supply a generic circuit contraction
bridge for this task. No such bridge is used in the positive theorem.

\section{Other completed bounds and prior-art comparison}
A separate completed hybrid compiler treats a 2-CNF $B$ together with
direct clauses $C$ having $H$ literal incidences. It manufactures the
source parity summary for $B$ alone and attaches physical variable ports
for $C$, obtaining $O(N\log\log N+H)$ Boolean wires. The key invariant
is equality of the solution sets, not just equality of scalar counts.
Combining this with a disjoint long-clause split gives raw AND/OR
circuits with arbitrary fan-in and $g$ gates the bound, for $N=n+g$,
\[
 (L+N)^{O(1)}
 2^{O\left(N\sqrt{\log\log N}/(\log N)^{1/4}\right)}.
\]
It permits dense wiring, but still pays for the gate variables.
The qualitative $g=O(n)$ consequence is already obtainable from fixed
width and a simple split, so it is not the principal advance claimed here.
An earlier uniform growing-width theorem and its independent audits are
also retained in the archive.

The targeted primary-source comparison supports a substantive conditional
extension, not a first-ever priority claim. Williams \cite{williams2018}
separates exact constant-degree polynomial counting from approximate
quadratic-system counting. Complex weighted counting dichotomies
\cite{liu} provide relevant hardness and tractable classes, but not a
reverse original-variable-preserving encoding eliminating arbitrary
factor density. Codsi--van de Wetering \cite{codsi} exploit random graph
structure for exact IQP improvements; their sparse-instance bound gives
savings from $2^n$, rather than this worst-case $2^{o(n)}$ guarantee.
Wahl--Strelchuk \cite{wahl} already give a genuine subexponential
simulation result under finite-range fixed-dimensional geometry.
De Colnet et al. \cite{decolnet} construct path representations and
evaluate them through rank-width. Small rank-width can cover instances
outside our density range, so neither method uniformly dominates the other.

A selective search of the pinned OpenAI release catalog checked the
entries on noncommutative formula identity testing, exact Fourier
transforms, finite-monoid expressions, one-tape space simulation, and
the governing counting result. Their stated contracts did not supply
a ready-made time-efficient generic-program compression step. This
is a search finding, not an impossibility claim or an exhaustive archive
review.

\section{Executed programs and preserved sources}\label{sec:executed}
All new finite validation reported here used exact integer or rational
arithmetic and passed. Distinct checkers were written independently.
The rows below count their actual executed fixtures; they are not timing
benchmarks or executions of the full asymptotic solver.

\begin{center}\small
\begin{tabular}{>{\raggedright\arraybackslash}p{0.43\linewidth}>{\raggedright\arraybackslash}p{0.48\linewidth}}
\toprule
Executed component & Verified scope\\\midrule
Principal phase construction & 30 truncation cases at orders $2,4,8$;
18 raw arithmetic programs; 6 multiple-output reductions; 168 IQP prefix identities.\\[4pt]
Independent phase audit & 72 fixtures; 336 character reconstructions;
insufficient-precision and single-character-histogram counterexamples.\\[4pt]
Root Boolean program compiler & 21 raw programs; 1,524 input assignments;
84 exact target-preimage counts; 9 cyclotomic phase cases;
93 IQP prefix identities; explicit rejection of a degree violation.\\[4pt]
General H/diagonal path compiler & 24 circuits, including arbitrary local
phase tables; 144 amplitude and 264 marginal comparisons against exact
state vectors; 2 inconsistent boundaries and 2 constant-wire table checks.\\[4pt]
Dense-grid obstruction & 29 exact spectral-rank certificates and 120
signed-weight cases.\\[4pt]
Hybrid binary/direct-clause compiler & 57 fixtures; 9,561 incidence subsets;
38,244 subset-index round trips; 1,192 boundary-moment identities.\\\bottomrule
\end{tabular}
\end{center}

The finite counters use exhaustive assignments, exact subset-union dynamic
programming, or exhaustive finite contractions. Their purpose is to
validate compilation, truncation, observables, and reconstruction.
There is no practical speed claim and therefore no claim of beating
equally informed deployed competitors. Implementing and optimizing the
complete replacement solver remains separate engineering work.

The supplied latest passive-network source, its full scientific audit,
and its receipt are preserved unchanged in the archive. Their reported
290/932-reading experiments are producer reports, not replayed network
experiments from this turn. The receipt's distinction between conditional
matrix concavity, derivative bounds, and physically acquired information
is retained. The older 510-reaction exact/noise-aware rank-one Gram source
is referenced by these documents, but its original proof and program
files were not recovered through the available selective routes. No
claim of preserving or replaying those missing bytes is made.

\section{Scientific boundary and compact continuation}
The established extension is a reusable representation principle:
\emph{bounded final coefficient height plus divisible local increments
can make reconstruction degree depend on original-input precision,
rather than on the number of operations}. It applies here because
fixed $2$-power phases have both properties, and the manufactured
summary cost tolerates incidence $n^{o(1)}$.

It does not cover arbitrary dense cubic polynomials with $\Theta(n^3)$
terms, unrestricted arithmetic degree, unrestricted Hadamard count,
arbitrary rational weights, or generic tensor networks. For density
$m=n^{1+\varepsilon}$ with fixed positive $\varepsilon$, the controlling
ratio $h/\log n$ does not vanish. Generic output support need not be
small or cheaply certifiable. The proof is conditional on the governing
manuscript, internally audited here, and has not received external peer
review.

The next substantive frontier is now exact: remove the growing-incidence
cost for dense cubic phases or for dense affine/quadratic preprocessing
without converting it into a large scope or many independent trace bits.
The signed complete-grid shortcut and output-normalization shortcut have
explicit obstructions in the accompanying proofs. A continuation should
start from the phase filtration and its exact bit costs, not from those
rejected shortcuts or from the earlier dense-CNF corollary. All actual
proofs and executed sources needed to resume are included; this statement
does not imply computation continues after the response ends.

\begin{thebibliography}{9}\small
\bibitem{governing} OpenAI. \emph{The Exponential Time Hypothesis is False}.
25 October 2026. Pinned commit
\texttt{7ef5c0cc5ca4b0e072b7a050fca09ab24bc6fa3a}.
\href{https://github.com/openai/math/blob/7ef5c0cc5ca4b0e072b7a050fca09ab24bc6fa3a/preprints/The-Exponential-Time-Hypothesis-is-False-October-25-2026/paper.pdf}{Governing PDF}.
Matched Git blob \texttt{0503566e95bc587a4076573a45132d97531b8ba6}.

\bibitem{williams2018} R. Ryan Williams.
\emph{Counting Solutions to Polynomial Systems via Reductions}. SOSA 2018.
\url{https://people.csail.mit.edu/rrw/counting-sosa.pdf}.

\bibitem{liu} Ying Liu.
\emph{The Exponential-Time Complexity of the Complex Weighted \#CSP}. 2022.
\url{https://arxiv.org/html/2202.02782v1}.

\bibitem{codsi} Julien Codsi and John van de Wetering.
\emph{Classically Simulating Quantum Supremacy IQP Circuits through a
Random Graph Approach}. Primary manuscript, 13 November 2024.
\url{https://ir.cwi.nl/pub/35965/35965.pdf}.

\bibitem{wahl} Thorsten B. Wahl and Sergii Strelchuk.
\emph{Simulating Quantum Circuits Using Efficient Tensor Network
Contraction Algorithms with Subexponential Upper Bound}.
Physical Review Letters 131, 180601 (2023).
\url{https://arxiv.org/abs/2208.01498}.

\bibitem{decolnet} Alexis de Colnet et al.
\emph{Quadratic Sums-of-Powers for Fixed-Parameter Tractable
Quantum-Circuit Simulation}. 28 May 2026.
\url{https://arxiv.org/html/2605.29944v1}.

\bibitem{n1077} Astra research knowledge. N1077 packet and its linked
governing tensor proof. Pinned commit
\texttt{e1e2272165c8548326c31bb11817c28f871d21e9}.
\href{https://github.com/harjassand/astra-research-knowledge/blob/e1e2272165c8548326c31bb11817c28f871d21e9/knowledge/current/packets/N1077.json}{Packet and source routes}.
\end{thebibliography}
\end{document}
